% appendix/D-reproducibility.tex
\chapter{ความสามารถในการทำซ้ำ}
\label{app:repro}

\section{สภาพแวดล้อมการรัน}

\begin{table}[H]
\centering
\caption{สภาพแวดล้อมการรันของผลการศึกษา}
\begin{tabularx}{\textwidth}{lX}
\toprule
\textbf{ส่วนประกอบ} & \textbf{เวอร์ชัน / รายละเอียด} \\
\midrule
ระบบปฏิบัติการ (local)        & macOS 15.2 (Apple M1, 8 GB unified memory) \\
ระบบปฏิบัติการ (cloud)        & Ubuntu 22.04 (Hugging Face Jobs container) \\
Python                        & 3.11.x \\
ผู้จัดการแพ็กเกจ              & uv 0.5.x \\
PyTorch                       & 2.5.x (CUDA 12.4 build บนคลาวด์) \\
transformers                  & 5.0.x \\
datasets                      & 3.0.x \\
scikit-learn                  & 1.5.x \\
LightGBM                      & 4.5.x \\
XGBoost                       & 2.1.x \\
SHAP                          & 0.46.x \\
LIME                          & 0.2.0.1 \\
MLflow                        & 2.18.x \\
GPU (cloud encoder FT)        & NVIDIA T4 (16 GB) / A10G (24 GB) \\
\bottomrule
\end{tabularx}
\end{table}

\section{คำสั่งที่ใช้ในการทำซ้ำผล}

\begin{lstlisting}[language=bash,caption={ขั้นตอนการ reproduce ทั้งหมด}]
# 1. คลังโค้ดและข้อมูล (ขอจากคณะผู้วิจัยหรืออาจารย์ที่ปรึกษา)
#    \authorone, \authortwo, \authorthree
#    Advisor: \advisor

# 2. ติดตั้ง environment
uv venv --python 3.11
uv pip install -e ".[dev]"

# 3. ใส่ API key (สร้างไฟล์ .env.local)
echo "YOUTUBE_API_KEY=YOUR_KEY_HERE" >> .env.local
echo "HUGGINGFACE_TOKEN=YOUR_HF_TOKEN" >> .env.local

# 4. เก็บข้อมูลและจัดเตรียม
make collect            # YouTube Data API v3 (1-2 hours)
make prepare            # clean + label + split + features

# 5. ฝึก baselines (local CPU)
make train-baselines    # LR / LightGBM / XGBoost (~10 min)
make train-hybrid       # MLP + GBM heads (~5 min)

# 6. ฝึก transformers (cloud GPU via HF Jobs)
make train-hf-all-detached
# poll: scripts/cloud/run_on_hf_jobs.py --poll-only --job-id <id>

# 7. ประกอบและประเมินผล
make train-stacking     # meta-LR + Platt calibration
make stacking-drop-one  # drop-one ablation
make eval               # ROC, McNemar, Cochran's Q, calibration
make explain            # SHAP / LIME / attention rollout

# 8. สร้างกราฟ
uv run python scripts/thesis_figures.py

# 9. คอมไพล์ thesis
cd docs/thesis
xelatex main && biber main && xelatex main && xelatex main
\end{lstlisting}

\section{Hash ของข้อมูลและ git SHA}

\begin{table}[H]
\centering
\caption{Hash อ้างอิงของผลการศึกษา}
\begin{tabularx}{\textwidth}{lX}
\toprule
\textbf{Artifact} & \textbf{SHA-256 (12 ตัวแรก) / git SHA} \\
\midrule
\texttt{data/processed/dataset\_with\_labels.parquet} & บันทึกใน MLflow tag \texttt{dataset\_hash} \\
git SHA ของผลในบทที่ \ref{ch:results}                  & \texttt{60fad4a} (พร้อมการเพิ่ม HF Jobs) \\
MLflow experiment                                      & \texttt{thesis-virality-thai} \\
HF Datasets (private, จัดเตรียมโดยคณะผู้วิจัย)         & input + per-model outputs \\
\bottomrule
\end{tabularx}
\end{table}

\section{เหตุการณ์และข้อจำกัดเชิงเทคนิค}

\subsection{Apple M1 MPS และ NaN logits}
ในระยะแรกของการทดลอง คณะผู้วิจัยพยายามฝึก PhayaThaiBERT บนเครื่อง Apple M1
ผ่าน Metal Performance Shaders (MPS) backend พบว่าการตั้ง \texttt{fp16=True}
หรือ \texttt{bf16=True} ส่งผลให้ logits เป็น NaN ตั้งแต่ epoch แรก คณะผู้วิจัย
แก้ไขด้วยการเพิ่มฟังก์ชัน \texttt{\_has\_cuda()} ใน
\texttt{src/models/transformer\_finetune.py:16-25} ที่ปิด mixed precision
อัตโนมัติเมื่อไม่ใช่ CUDA

\subsection{Kaggle Free GPU Pool}
คิว Kaggle Free GPU pool ในช่วงดำเนินงานหนาแน่นต่อเนื่องเกิน 6 ชั่วโมง คณะ
ผู้วิจัยจึงตัดสินใจย้ายการฝึกไปยัง Hugging Face Jobs ซึ่งจองได้ทันที ค่าใช้จ่าย
รวมประมาณ \$0.50

\subsection{Train\_hybrid Silent GBM Skip on MPS}
สคริปต์ \texttt{scripts/train\_hybrid.py} แสดงพฤติกรรมที่ข้าม training ของ
GBM head อย่างเงียบ ๆ หลังจากฝึก MLP บน MPS เนื่องจาก backend resource
contention คณะผู้วิจัยใช้ workaround โดยรัน MLP และ GBM ในสคริปต์แยก
รายละเอียดอยู่ใน \texttt{.claude/memory/decisions.md}

\section{การ Audit ทางสถิติ}

ผลทุกตัวเลขในบทที่ \ref{ch:results} ถูกตรวจสอบความสอดคล้องระหว่าง 3 แหล่ง
ดังนี้

\begin{enumerate}[leftmargin=2em]
  \item \textbf{MLflow run UI} — point estimate และค่า metric ทุก epoch
  \item \textbf{CSV ใน \texttt{reports/tables/}} — point estimate + bootstrap
        CI 95\% สำหรับใช้พิมพ์ในตาราง LaTeX
  \item \textbf{Parquet ใน \texttt{reports/artifacts/predictions/}} —
        \texttt{(video\_id, split, y\_true, y\_proba)} สำหรับ McNemar /
        Cochran / re-run bootstrap
\end{enumerate}

หากพบความไม่สอดคล้องระหว่างสามแหล่ง parquet predictions เป็น source of
truth ที่ recompute ได้ทันทีโดยไม่ต้อง re-train
